%%%%%%%%%%%%%%%%%%%%%%%%%%%%%%%%%%%%%%%%%%%%%%%%%%%%%%%%%%%%%%%%%%%%%
%  RISE_main.tex
%  Research in Science Education (Springer) 투고용 LaTeX 템플릿
%  ── 공식 Springer Nature 템플릿(sn-jnl.cls, v3.x) 기반 ──
%
%  ▶ 이 템플릿은 무엇인가
%    - Springer 공식 LaTeX 클래스 sn-jnl.cls 를 사용합니다(폴더에 함께 넣어 둔
%      sn-jnl.cls / sn-apacite.bst 사용). 게재 확정 후 production 단계와 가장 잘
%      맞습니다.
%    - 본문은 sub/en-1..5-*.tex 섹션 파일로 분리되어 있으며, 예시 내용 없이 각
%      절에 무엇을 쓰는지 안내만 넣어 두었습니다. 새 원고는 이 파일들을 채워
%      작성하면 됩니다.
%    - biblatex 인용 명령(\parencite/\textcite)을 그대로 쓸 수 있도록 아래에
%      natbib 호환 shim 을 넣어 두었습니다(sn-jnl 은 natbib 기반).
%
%  ▶ 참고문헌: sn-apa 옵션 → sn-apacite.bst(APA) 자동 지정. BibTeX 사용.
%  ▶ 컴파일(pdfLaTeX + BibTeX):
%      pdflatex RISE_main → bibtex RISE_main → pdflatex RISE_main → pdflatex RISE_main
%      (또는 제공된 RISE_main_build.cmd / RISE_main_build.sh 사용)
%
%  ▶ 클래스 옵션
%      pdflatex : pdfLaTeX 엔진      / sn-apa  : APA 참고문헌
%      referee  : 2줄 간격(심사 투고용, 기본값) / lineno : 줄번호(필요 시 추가)
%      게재 확정본(최종본)을 조판할 때는 referee 를 빼면 1줄 간격 저널 레이아웃이
%      됩니다.
%      ※ 이중 익명(double-blind) 투고 시에는 아래 \author*/\affil* 블록의 식별
%        정보를 제거한 익명본을 별도로 만드세요(sn-jnl 은 자동 익명화하지 않음).
%
%  ▶ 투고 전 확인: Research in Science Education(Springer) 공식 Instructions for
%    Authors 에서 원고 유형, 초록/키워드 요구사항, 그림/표 규격, 윤리·이해상충
%    선언 항목을 반드시 최신본으로 확인하세요.
%%%%%%%%%%%%%%%%%%%%%%%%%%%%%%%%%%%%%%%%%%%%%%%%%%%%%%%%%%%%%%%%%%%%%
\documentclass[pdflatex,sn-apa,referee]{sn-jnl}

% --- sn-jnl.cls 호환 패치 -------------------------------------------------
% 공식 sn-jnl.cls 는 begindocument 단계에서 \SetFootnoteHook /
% \DeclareNewFootnote(= manyfoot 명령)를 사용하는데, 최신 TeX Live 의 footmisc 가
% manyfoot 를 자동으로 끌어오지 않아 미정의가 될 수 있습니다. manyfoot 를 직접
% 로드하여 해결합니다. (향후 Springer 가 이 의존성을 cls 내부에서 처리하면 제거 가능)
\usepackage{manyfoot}

% --- 본문(en-*.tex)에서 흔히 쓰는 패키지 (중복 로드는 무해) ---------------
\usepackage{graphicx}
\usepackage{tikz}\usetikzlibrary{arrows.meta}
\usepackage{booktabs}
\usepackage{tabularx}
\usepackage{array}
\usepackage{multirow}
\usepackage{amsmath,amssymb}
% sn-jnl.cls 는 자체 \orcidlogo(Orcidlogo.eps + shell-escape 필요)를 정의하므로
% orcidlink 와 충돌합니다. \orcidlogo 를 풀어 orcidlink(tikz 기반, eps 불필요)가
% 쓰이게 합니다. 이후 \orcidlink{0000-...} 로 ORCID 아이콘+링크를 출력합니다.
\let\orcidlogo\relax
\usepackage{orcidlink}   % 저자 ORCID 링크(아이콘, eps 불필요)

% --- biblatex → natbib 인용 명령 호환 shim --------------------------------
% 본문에서 \parencite/\textcite(biblatex 문법)를 써도 되도록 natbib 의
% \citep/\citet 로 매핑합니다. (sn-jnl 은 natbib 기반)
\newcommand{\parencite}[1]{\citep{#1}}
\newcommand{\textcite}[1]{\citet{#1}}

% --- 표 헬퍼 --------------------------------------------------------------
% \TableNote{...} : 표 아래 Note. 주석 / L 열 : 좌측 정렬 가변폭(tabularx)
\newcommand{\TableNote}[1]{\par\vspace{2pt}\noindent\parbox{\linewidth}{\footnotesize\raggedright\textit{Note.}\ #1}}
\newcolumntype{L}{>{\raggedright\arraybackslash}X}

% 게재본은 절 제목에 번호를 달지 않습니다(sn-jnl 제공 명령, secnumdepth=0).
% 번호를 달고 싶으면 아래 줄을 주석 처리하세요.
\unnumbered

% 참고문헌은 단일 간격으로 조판합니다(referee 더블스페이싱이 항목 간격을
% 들쭉날쭉하게 만드는 것을 방지 — 본문은 referee 시 더블 간격 유지).
\usepackage{etoolbox}
\AtBeginEnvironment{thebibliography}{\singlespacing}

% --- 이중 익명(double-blind) 토글 -----------------------------------------
% RISE(Research in Science Education)는 이중 익명 심사입니다. 익명본을 만들 때는
% 빌드에서  \def\ANON{}%%%%%%%%%%%%%%%%%%%%%%%%%%%%%%%%%%%%%%%%%%%%%%%%%%%%%%%%%%%%%%%%%%%%%
%  RISE_main.tex
%  Research in Science Education (Springer) 투고용 LaTeX 템플릿
%  ── 공식 Springer Nature 템플릿(sn-jnl.cls, v3.x) 기반 ──
%
%  ▶ 이 템플릿은 무엇인가
%    - Springer 공식 LaTeX 클래스 sn-jnl.cls 를 사용합니다(폴더에 함께 넣어 둔
%      sn-jnl.cls / sn-apacite.bst 사용). 게재 확정 후 production 단계와 가장 잘
%      맞습니다.
%    - 본문은 sub/en-1..5-*.tex 섹션 파일로 분리되어 있으며, 예시 내용 없이 각
%      절에 무엇을 쓰는지 안내만 넣어 두었습니다. 새 원고는 이 파일들을 채워
%      작성하면 됩니다.
%    - biblatex 인용 명령(\parencite/\textcite)을 그대로 쓸 수 있도록 아래에
%      natbib 호환 shim 을 넣어 두었습니다(sn-jnl 은 natbib 기반).
%
%  ▶ 참고문헌: sn-apa 옵션 → sn-apacite.bst(APA) 자동 지정. BibTeX 사용.
%  ▶ 컴파일(pdfLaTeX + BibTeX):
%      pdflatex RISE_main → bibtex RISE_main → pdflatex RISE_main → pdflatex RISE_main
%      (또는 제공된 RISE_main_build.cmd / RISE_main_build.sh 사용)
%
%  ▶ 클래스 옵션
%      pdflatex : pdfLaTeX 엔진      / sn-apa  : APA 참고문헌
%      referee  : 2줄 간격(심사 투고용, 기본값) / lineno : 줄번호(필요 시 추가)
%      게재 확정본(최종본)을 조판할 때는 referee 를 빼면 1줄 간격 저널 레이아웃이
%      됩니다.
%      ※ 이중 익명(double-blind) 투고 시에는 아래 \author*/\affil* 블록의 식별
%        정보를 제거한 익명본을 별도로 만드세요(sn-jnl 은 자동 익명화하지 않음).
%
%  ▶ 투고 전 확인: Research in Science Education(Springer) 공식 Instructions for
%    Authors 에서 원고 유형, 초록/키워드 요구사항, 그림/표 규격, 윤리·이해상충
%    선언 항목을 반드시 최신본으로 확인하세요.
%%%%%%%%%%%%%%%%%%%%%%%%%%%%%%%%%%%%%%%%%%%%%%%%%%%%%%%%%%%%%%%%%%%%%
\documentclass[pdflatex,sn-apa,referee]{sn-jnl}

% --- sn-jnl.cls 호환 패치 -------------------------------------------------
% 공식 sn-jnl.cls 는 begindocument 단계에서 \SetFootnoteHook /
% \DeclareNewFootnote(= manyfoot 명령)를 사용하는데, 최신 TeX Live 의 footmisc 가
% manyfoot 를 자동으로 끌어오지 않아 미정의가 될 수 있습니다. manyfoot 를 직접
% 로드하여 해결합니다. (향후 Springer 가 이 의존성을 cls 내부에서 처리하면 제거 가능)
\usepackage{manyfoot}

% --- 본문(en-*.tex)에서 흔히 쓰는 패키지 (중복 로드는 무해) ---------------
\usepackage{graphicx}
\usepackage{tikz}\usetikzlibrary{arrows.meta}
\usepackage{booktabs}
\usepackage{tabularx}
\usepackage{array}
\usepackage{multirow}
\usepackage{amsmath,amssymb}
% sn-jnl.cls 는 자체 \orcidlogo(Orcidlogo.eps + shell-escape 필요)를 정의하므로
% orcidlink 와 충돌합니다. \orcidlogo 를 풀어 orcidlink(tikz 기반, eps 불필요)가
% 쓰이게 합니다. 이후 \orcidlink{0000-...} 로 ORCID 아이콘+링크를 출력합니다.
\let\orcidlogo\relax
\usepackage{orcidlink}   % 저자 ORCID 링크(아이콘, eps 불필요)

% --- biblatex → natbib 인용 명령 호환 shim --------------------------------
% 본문에서 \parencite/\textcite(biblatex 문법)를 써도 되도록 natbib 의
% \citep/\citet 로 매핑합니다. (sn-jnl 은 natbib 기반)
\newcommand{\parencite}[1]{\citep{#1}}
\newcommand{\textcite}[1]{\citet{#1}}

% --- 표 헬퍼 --------------------------------------------------------------
% \TableNote{...} : 표 아래 Note. 주석 / L 열 : 좌측 정렬 가변폭(tabularx)
\newcommand{\TableNote}[1]{\par\vspace{2pt}\noindent\parbox{\linewidth}{\footnotesize\raggedright\textit{Note.}\ #1}}
\newcolumntype{L}{>{\raggedright\arraybackslash}X}

% 게재본은 절 제목에 번호를 달지 않습니다(sn-jnl 제공 명령, secnumdepth=0).
% 번호를 달고 싶으면 아래 줄을 주석 처리하세요.
\unnumbered

% 참고문헌은 단일 간격으로 조판합니다(referee 더블스페이싱이 항목 간격을
% 들쭉날쭉하게 만드는 것을 방지 — 본문은 referee 시 더블 간격 유지).
\usepackage{etoolbox}
\AtBeginEnvironment{thebibliography}{\singlespacing}

% --- 이중 익명(double-blind) 토글 -----------------------------------------
% RISE(Research in Science Education)는 이중 익명 심사입니다. 익명본을 만들 때는
% 빌드에서  \def\ANON{}%%%%%%%%%%%%%%%%%%%%%%%%%%%%%%%%%%%%%%%%%%%%%%%%%%%%%%%%%%%%%%%%%%%%%
%  RISE_main.tex
%  Research in Science Education (Springer) 투고용 LaTeX 템플릿
%  ── 공식 Springer Nature 템플릿(sn-jnl.cls, v3.x) 기반 ──
%
%  ▶ 이 템플릿은 무엇인가
%    - Springer 공식 LaTeX 클래스 sn-jnl.cls 를 사용합니다(폴더에 함께 넣어 둔
%      sn-jnl.cls / sn-apacite.bst 사용). 게재 확정 후 production 단계와 가장 잘
%      맞습니다.
%    - 본문은 sub/en-1..5-*.tex 섹션 파일로 분리되어 있으며, 예시 내용 없이 각
%      절에 무엇을 쓰는지 안내만 넣어 두었습니다. 새 원고는 이 파일들을 채워
%      작성하면 됩니다.
%    - biblatex 인용 명령(\parencite/\textcite)을 그대로 쓸 수 있도록 아래에
%      natbib 호환 shim 을 넣어 두었습니다(sn-jnl 은 natbib 기반).
%
%  ▶ 참고문헌: sn-apa 옵션 → sn-apacite.bst(APA) 자동 지정. BibTeX 사용.
%  ▶ 컴파일(pdfLaTeX + BibTeX):
%      pdflatex RISE_main → bibtex RISE_main → pdflatex RISE_main → pdflatex RISE_main
%      (또는 제공된 RISE_main_build.cmd / RISE_main_build.sh 사용)
%
%  ▶ 클래스 옵션
%      pdflatex : pdfLaTeX 엔진      / sn-apa  : APA 참고문헌
%      referee  : 2줄 간격(심사 투고용, 기본값) / lineno : 줄번호(필요 시 추가)
%      게재 확정본(최종본)을 조판할 때는 referee 를 빼면 1줄 간격 저널 레이아웃이
%      됩니다.
%      ※ 이중 익명(double-blind) 투고 시에는 아래 \author*/\affil* 블록의 식별
%        정보를 제거한 익명본을 별도로 만드세요(sn-jnl 은 자동 익명화하지 않음).
%
%  ▶ 투고 전 확인: Research in Science Education(Springer) 공식 Instructions for
%    Authors 에서 원고 유형, 초록/키워드 요구사항, 그림/표 규격, 윤리·이해상충
%    선언 항목을 반드시 최신본으로 확인하세요.
%%%%%%%%%%%%%%%%%%%%%%%%%%%%%%%%%%%%%%%%%%%%%%%%%%%%%%%%%%%%%%%%%%%%%
\documentclass[pdflatex,sn-apa,referee]{sn-jnl}

% --- sn-jnl.cls 호환 패치 -------------------------------------------------
% 공식 sn-jnl.cls 는 begindocument 단계에서 \SetFootnoteHook /
% \DeclareNewFootnote(= manyfoot 명령)를 사용하는데, 최신 TeX Live 의 footmisc 가
% manyfoot 를 자동으로 끌어오지 않아 미정의가 될 수 있습니다. manyfoot 를 직접
% 로드하여 해결합니다. (향후 Springer 가 이 의존성을 cls 내부에서 처리하면 제거 가능)
\usepackage{manyfoot}

% --- 본문(en-*.tex)에서 흔히 쓰는 패키지 (중복 로드는 무해) ---------------
\usepackage{graphicx}
\usepackage{tikz}\usetikzlibrary{arrows.meta}
\usepackage{booktabs}
\usepackage{tabularx}
\usepackage{array}
\usepackage{multirow}
\usepackage{amsmath,amssymb}
% sn-jnl.cls 는 자체 \orcidlogo(Orcidlogo.eps + shell-escape 필요)를 정의하므로
% orcidlink 와 충돌합니다. \orcidlogo 를 풀어 orcidlink(tikz 기반, eps 불필요)가
% 쓰이게 합니다. 이후 \orcidlink{0000-...} 로 ORCID 아이콘+링크를 출력합니다.
\let\orcidlogo\relax
\usepackage{orcidlink}   % 저자 ORCID 링크(아이콘, eps 불필요)

% --- biblatex → natbib 인용 명령 호환 shim --------------------------------
% 본문에서 \parencite/\textcite(biblatex 문법)를 써도 되도록 natbib 의
% \citep/\citet 로 매핑합니다. (sn-jnl 은 natbib 기반)
\newcommand{\parencite}[1]{\citep{#1}}
\newcommand{\textcite}[1]{\citet{#1}}

% --- 표 헬퍼 --------------------------------------------------------------
% \TableNote{...} : 표 아래 Note. 주석 / L 열 : 좌측 정렬 가변폭(tabularx)
\newcommand{\TableNote}[1]{\par\vspace{2pt}\noindent\parbox{\linewidth}{\footnotesize\raggedright\textit{Note.}\ #1}}
\newcolumntype{L}{>{\raggedright\arraybackslash}X}

% 게재본은 절 제목에 번호를 달지 않습니다(sn-jnl 제공 명령, secnumdepth=0).
% 번호를 달고 싶으면 아래 줄을 주석 처리하세요.
\unnumbered

% 참고문헌은 단일 간격으로 조판합니다(referee 더블스페이싱이 항목 간격을
% 들쭉날쭉하게 만드는 것을 방지 — 본문은 referee 시 더블 간격 유지).
\usepackage{etoolbox}
\AtBeginEnvironment{thebibliography}{\singlespacing}

% --- 이중 익명(double-blind) 토글 -----------------------------------------
% RISE(Research in Science Education)는 이중 익명 심사입니다. 익명본을 만들 때는
% 빌드에서  \def\ANON{}%%%%%%%%%%%%%%%%%%%%%%%%%%%%%%%%%%%%%%%%%%%%%%%%%%%%%%%%%%%%%%%%%%%%%
%  RISE_main.tex
%  Research in Science Education (Springer) 투고용 LaTeX 템플릿
%  ── 공식 Springer Nature 템플릿(sn-jnl.cls, v3.x) 기반 ──
%
%  ▶ 이 템플릿은 무엇인가
%    - Springer 공식 LaTeX 클래스 sn-jnl.cls 를 사용합니다(폴더에 함께 넣어 둔
%      sn-jnl.cls / sn-apacite.bst 사용). 게재 확정 후 production 단계와 가장 잘
%      맞습니다.
%    - 본문은 sub/en-1..5-*.tex 섹션 파일로 분리되어 있으며, 예시 내용 없이 각
%      절에 무엇을 쓰는지 안내만 넣어 두었습니다. 새 원고는 이 파일들을 채워
%      작성하면 됩니다.
%    - biblatex 인용 명령(\parencite/\textcite)을 그대로 쓸 수 있도록 아래에
%      natbib 호환 shim 을 넣어 두었습니다(sn-jnl 은 natbib 기반).
%
%  ▶ 참고문헌: sn-apa 옵션 → sn-apacite.bst(APA) 자동 지정. BibTeX 사용.
%  ▶ 컴파일(pdfLaTeX + BibTeX):
%      pdflatex RISE_main → bibtex RISE_main → pdflatex RISE_main → pdflatex RISE_main
%      (또는 제공된 RISE_main_build.cmd / RISE_main_build.sh 사용)
%
%  ▶ 클래스 옵션
%      pdflatex : pdfLaTeX 엔진      / sn-apa  : APA 참고문헌
%      referee  : 2줄 간격(심사 투고용, 기본값) / lineno : 줄번호(필요 시 추가)
%      게재 확정본(최종본)을 조판할 때는 referee 를 빼면 1줄 간격 저널 레이아웃이
%      됩니다.
%      ※ 이중 익명(double-blind) 투고 시에는 아래 \author*/\affil* 블록의 식별
%        정보를 제거한 익명본을 별도로 만드세요(sn-jnl 은 자동 익명화하지 않음).
%
%  ▶ 투고 전 확인: Research in Science Education(Springer) 공식 Instructions for
%    Authors 에서 원고 유형, 초록/키워드 요구사항, 그림/표 규격, 윤리·이해상충
%    선언 항목을 반드시 최신본으로 확인하세요.
%%%%%%%%%%%%%%%%%%%%%%%%%%%%%%%%%%%%%%%%%%%%%%%%%%%%%%%%%%%%%%%%%%%%%
\documentclass[pdflatex,sn-apa,referee]{sn-jnl}

% --- sn-jnl.cls 호환 패치 -------------------------------------------------
% 공식 sn-jnl.cls 는 begindocument 단계에서 \SetFootnoteHook /
% \DeclareNewFootnote(= manyfoot 명령)를 사용하는데, 최신 TeX Live 의 footmisc 가
% manyfoot 를 자동으로 끌어오지 않아 미정의가 될 수 있습니다. manyfoot 를 직접
% 로드하여 해결합니다. (향후 Springer 가 이 의존성을 cls 내부에서 처리하면 제거 가능)
\usepackage{manyfoot}

% --- 본문(en-*.tex)에서 흔히 쓰는 패키지 (중복 로드는 무해) ---------------
\usepackage{graphicx}
\usepackage{tikz}\usetikzlibrary{arrows.meta}
\usepackage{booktabs}
\usepackage{tabularx}
\usepackage{array}
\usepackage{multirow}
\usepackage{amsmath,amssymb}
% sn-jnl.cls 는 자체 \orcidlogo(Orcidlogo.eps + shell-escape 필요)를 정의하므로
% orcidlink 와 충돌합니다. \orcidlogo 를 풀어 orcidlink(tikz 기반, eps 불필요)가
% 쓰이게 합니다. 이후 \orcidlink{0000-...} 로 ORCID 아이콘+링크를 출력합니다.
\let\orcidlogo\relax
\usepackage{orcidlink}   % 저자 ORCID 링크(아이콘, eps 불필요)

% --- biblatex → natbib 인용 명령 호환 shim --------------------------------
% 본문에서 \parencite/\textcite(biblatex 문법)를 써도 되도록 natbib 의
% \citep/\citet 로 매핑합니다. (sn-jnl 은 natbib 기반)
\newcommand{\parencite}[1]{\citep{#1}}
\newcommand{\textcite}[1]{\citet{#1}}

% --- 표 헬퍼 --------------------------------------------------------------
% \TableNote{...} : 표 아래 Note. 주석 / L 열 : 좌측 정렬 가변폭(tabularx)
\newcommand{\TableNote}[1]{\par\vspace{2pt}\noindent\parbox{\linewidth}{\footnotesize\raggedright\textit{Note.}\ #1}}
\newcolumntype{L}{>{\raggedright\arraybackslash}X}

% 게재본은 절 제목에 번호를 달지 않습니다(sn-jnl 제공 명령, secnumdepth=0).
% 번호를 달고 싶으면 아래 줄을 주석 처리하세요.
\unnumbered

% 참고문헌은 단일 간격으로 조판합니다(referee 더블스페이싱이 항목 간격을
% 들쭉날쭉하게 만드는 것을 방지 — 본문은 referee 시 더블 간격 유지).
\usepackage{etoolbox}
\AtBeginEnvironment{thebibliography}{\singlespacing}

% --- 이중 익명(double-blind) 토글 -----------------------------------------
% RISE(Research in Science Education)는 이중 익명 심사입니다. 익명본을 만들 때는
% 빌드에서  \def\ANON{}\input{RISE_main.tex}  로 컴파일합니다(제공된 빌드 스크립트의
% `anon` 명령이 자동 처리하여 RISE_main_anon.pdf 를 만듭니다). \ANON 이 정의되면
% 저자·소속·이메일·ORCID·사사(Funding)가 익명화됩니다. 저자 식별 정보는 익명 원고와
% 별도로 title-page.tex(투고용 표지)로 제출합니다.
\newif\ifRISEanon
\ifdefined\ANON\RISEanontrue\fi

% =========================================================================
% Front matter (sn-jnl 양식) — 아래 제목/저자/초록/키워드를 새 원고에 맞게 수정
% =========================================================================
\title[Short Running Title]{Full Article Title Goes Here:
Subtitle If Any}

% 저자 순서대로 \author 로 나열하고, 교신저자는 \author* 로 표시합니다.
% \ifRISEanon 이 참이면(익명본) 저자·소속·이메일이 익명 처리됩니다.
\ifRISEanon
  \author*[1]{\fnm{}\sur{Author name withheld for double-blind review}}
  \affil[1]{\orgname{Affiliation withheld for double-blind review}}
\else
  \author[1]{\fnm{First} \sur{Author}\,\orcidlink{0000-0000-0000-0000}}
  \author*[1]{\fnm{Corresponding} \sur{Author}\,\orcidlink{0000-0000-0000-0000}}\email{corresponding.author@example.edu}
  \author[2]{\fnm{Third} \sur{Author}}
  \affil[1]{\orgname{Your Institution}, \orgaddress{\city{City}, \country{Country}}}
  \affil[2]{\orgname{Second Institution}, \orgaddress{\city{City}, \country{Country}}}
\fi

% 초록: Research in Science Education 은 보통 구조화되지 않은(unstructured) 단일
% 문단 초록을 사용합니다. 배경/문제의식 → 목적·연구질문 → 대상·방법 → 핵심 결과
% → 과학교육 연구·실천에의 함의를 한 문단으로 압축합니다. 최신 저자 안내에서 분량
% 제한을 확인하세요. (이 클래스는 pdfLaTeX 로 컴파일되므로 본문은 영문으로 작성합니다.)
\abstract{Write the abstract here as a single unstructured paragraph. State the
background and problem, the purpose and research questions, the participants and
methods, the key findings, and the implications for science-education research
and practice. Check the journal's current author guidelines for length limits.}

\keywords{keyword one, keyword two, keyword three, keyword four}

\begin{document}

\maketitle

% =========================================================================
% 본문 — 각 절 파일은 sub/ 폴더에 있습니다(예시 내용 없이 안내만 포함).
% =========================================================================
\input{./sub/en-1-introduction.tex}   % 1. Introduction
\input{./sub/en-2-framework.tex}      % 2. Theoretical Framework
\input{./sub/en-3-methods.tex}        % 3. Methods
\input{./sub/en-4-results.tex}        % 4. Results
\input{./sub/en-5-discussion.tex}     % 5. Discussion and Conclusion

% =========================================================================
% Back matter
% =========================================================================
\backmatter

% Declarations: 아래 각 항목을 원고에 맞게 영문으로 작성합니다.
%   Funding(연구비) / Competing interests(이해상충) / Ethics approval(IRB 승인) /
%   Informed consent(연구참여 동의, 미성년자는 보호자 동의) /
%   Data availability(데이터 공개) / Generative AI use(생성형 AI 사용 범위)
\bmhead{Declarations}
\begin{itemize}
  \item \textbf{Funding.} State funding sources, or that no funding was received.
  \item \textbf{Competing interests.} State any competing interests.
  \item \textbf{Ethics approval.} State IRB/ethics approval information.
  \item \textbf{Informed consent.} State that informed consent (and guardian
        consent for minors) was obtained.
  \item \textbf{Data availability.} State how the data can be accessed or requested.
  \item \textbf{Generative AI use.} Disclose whether and how generative AI was used.
\end{itemize}

% 참고문헌(BibTeX, APA via sn-apacite). \bibliographystyle 은 sn-jnl 이 자동 지정.
\bibliography{sub/references-sn}

\end{document}
%%%%%%%%%%%%%%%%%%%%%%%%%%%%%%%%%%%%%%%%%%%%%%%%%%%%%%%%%%%%%%%%%%%%%
%  END RISE_main.tex
%%%%%%%%%%%%%%%%%%%%%%%%%%%%%%%%%%%%%%%%%%%%%%%%%%%%%%%%%%%%%%%%%%%%%
  로 컴파일합니다(제공된 빌드 스크립트의
% `anon` 명령이 자동 처리하여 RISE_main_anon.pdf 를 만듭니다). \ANON 이 정의되면
% 저자·소속·이메일·ORCID·사사(Funding)가 익명화됩니다. 저자 식별 정보는 익명 원고와
% 별도로 title-page.tex(투고용 표지)로 제출합니다.
\newif\ifRISEanon
\ifdefined\ANON\RISEanontrue\fi

% =========================================================================
% Front matter (sn-jnl 양식) — 아래 제목/저자/초록/키워드를 새 원고에 맞게 수정
% =========================================================================
\title[Short Running Title]{Full Article Title Goes Here:
Subtitle If Any}

% 저자 순서대로 \author 로 나열하고, 교신저자는 \author* 로 표시합니다.
% \ifRISEanon 이 참이면(익명본) 저자·소속·이메일이 익명 처리됩니다.
\ifRISEanon
  \author*[1]{\fnm{}\sur{Author name withheld for double-blind review}}
  \affil[1]{\orgname{Affiliation withheld for double-blind review}}
\else
  \author[1]{\fnm{First} \sur{Author}\,\orcidlink{0000-0000-0000-0000}}
  \author*[1]{\fnm{Corresponding} \sur{Author}\,\orcidlink{0000-0000-0000-0000}}\email{corresponding.author@example.edu}
  \author[2]{\fnm{Third} \sur{Author}}
  \affil[1]{\orgname{Your Institution}, \orgaddress{\city{City}, \country{Country}}}
  \affil[2]{\orgname{Second Institution}, \orgaddress{\city{City}, \country{Country}}}
\fi

% 초록: Research in Science Education 은 보통 구조화되지 않은(unstructured) 단일
% 문단 초록을 사용합니다. 배경/문제의식 → 목적·연구질문 → 대상·방법 → 핵심 결과
% → 과학교육 연구·실천에의 함의를 한 문단으로 압축합니다. 최신 저자 안내에서 분량
% 제한을 확인하세요. (이 클래스는 pdfLaTeX 로 컴파일되므로 본문은 영문으로 작성합니다.)
\abstract{Write the abstract here as a single unstructured paragraph. State the
background and problem, the purpose and research questions, the participants and
methods, the key findings, and the implications for science-education research
and practice. Check the journal's current author guidelines for length limits.}

\keywords{keyword one, keyword two, keyword three, keyword four}

\begin{document}

\maketitle

% =========================================================================
% 본문 — 각 절 파일은 sub/ 폴더에 있습니다(예시 내용 없이 안내만 포함).
% =========================================================================
%%%%%%%%%%%%%%%%%%%%%%%%%%%%%%%%%%%%%%%%%%%%%%%%%%%%%%%%%%%%%%%%%%%%%
%  sub/en-1-introduction.tex  —  1. Introduction
%%%%%%%%%%%%%%%%%%%%%%%%%%%%%%%%%%%%%%%%%%%%%%%%%%%%%%%%%%%%%%%%%%%%%
\section{Introduction}\label{sec:introduction}

% 서론에서는 연구가 다루는 과학교육 문제와 그 중요성을 제시하고, 선행연구를
% 비판적으로 검토하여 연구 공백(gap)을 드러낸 뒤, 본 연구의 목적과 연구 질문을
% 명확히 밝힙니다. Research in Science Education 은 이론적 기여와 실증적 근거의
% 연결을 중시하므로, 연구 질문이 어떤 이론적·실천적 함의를 겨냥하는지 초반에
% 분명히 하세요.

Write the introduction here. State the science-education problem, review prior
work to expose the gap, and end with explicit research questions.

% 인용은 \parencite{key} (괄호 인용) 또는 \textcite{key} (본문 인용)로 씁니다.
% 예: Prior work has examined this issue \parencite{example2024}, and
% \textcite{example2024} argued that \dots
   % 1. Introduction
%%%%%%%%%%%%%%%%%%%%%%%%%%%%%%%%%%%%%%%%%%%%%%%%%%%%%%%%%%%%%%%%%%%%%
%  sub/en-2-framework.tex  —  2. Theoretical Framework
%%%%%%%%%%%%%%%%%%%%%%%%%%%%%%%%%%%%%%%%%%%%%%%%%%%%%%%%%%%%%%%%%%%%%
\section{Theoretical Framework}\label{sec:framework}

% 연구를 뒷받침하는 이론적 틀, 핵심 개념, 분석 렌즈를 제시합니다. 연구 질문과
% 방법(코딩 체계, 루브릭 등)을 이 틀에 연결하면 이후 결과 해석의 근거가 됩니다.

Present the theoretical framework, key constructs, and analytic lens that ground
the study and connect the research questions to the methods.
      % 2. Theoretical Framework
%%%%%%%%%%%%%%%%%%%%%%%%%%%%%%%%%%%%%%%%%%%%%%%%%%%%%%%%%%%%%%%%%%%%%
%  sub/en-3-methods.tex  —  3. Methods
%%%%%%%%%%%%%%%%%%%%%%%%%%%%%%%%%%%%%%%%%%%%%%%%%%%%%%%%%%%%%%%%%%%%%
\section{Methods}\label{sec:methods}

% 연구 설계, 맥락과 참여자, 자료 수집, 분석 절차(코딩 체계·신뢰도·타당도)를
% 재현 가능하도록 기술합니다. 연구윤리(IRB) 승인과 동의 절차도 밝힙니다.

Describe the research design, context and participants, data sources, and the
analysis procedure (coding scheme, reliability, validity) in enough detail to be
reproducible.

% 표 예시 (필요 시 사용):
% \begin{table}[htbp]
%   \centering
%   \caption{Participant summary}\label{tab:participants}
%   \begin{tabularx}{\linewidth}{@{}lL@{}}
%     \toprule
%     Group & Description \\
%     \midrule
%     A & \dots \\
%     \bottomrule
%   \end{tabularx}
%   \TableNote{Add a table note here if needed.}
% \end{table}
        % 3. Methods
%%%%%%%%%%%%%%%%%%%%%%%%%%%%%%%%%%%%%%%%%%%%%%%%%%%%%%%%%%%%%%%%%%%%%
%  sub/en-4-results.tex  —  4. Results
%%%%%%%%%%%%%%%%%%%%%%%%%%%%%%%%%%%%%%%%%%%%%%%%%%%%%%%%%%%%%%%%%%%%%
\section{Results}\label{sec:results}

% 연구 질문 순서에 따라 결과를 제시합니다. 표와 그림은 본문에서 먼저 언급한 뒤
% 배치하고, 모든 표/그림에 \label 을 붙여 \ref/\cref 로 참조합니다. 해석은
% 최소화하고 근거(수치·인용·사례)를 중심으로 서술합니다.

Report the findings organized by research question. Reference every table and
figure in the text before it appears.

% 그림 예시 (필요 시 사용):
% \begin{figure}[htbp]
%   \centering
%   \includegraphics[width=\linewidth]{images/your-figure}
%   \caption{Describe the figure.}\label{fig:example}
% \end{figure}
        % 4. Results
%%%%%%%%%%%%%%%%%%%%%%%%%%%%%%%%%%%%%%%%%%%%%%%%%%%%%%%%%%%%%%%%%%%%%
%  sub/en-5-discussion.tex  —  5. Discussion and Conclusion
%%%%%%%%%%%%%%%%%%%%%%%%%%%%%%%%%%%%%%%%%%%%%%%%%%%%%%%%%%%%%%%%%%%%%
\section{Discussion and Conclusion}\label{sec:discussion}

% 핵심 결과를 연구 질문·이론 틀과 연결해 해석하고, 선행연구와의 관계(확장/반박),
% 과학교육 연구·실천에의 함의, 연구의 한계와 후속 연구 방향을 제시합니다.
% Research in Science Education 은 이론적·교육적 기여를 분명히 드러내는 논의를
% 중시합니다.

Interpret the findings in relation to the research questions and framework,
connect to prior work, state implications for science-education research and
practice, and acknowledge limitations and directions for future work.
     % 5. Discussion and Conclusion

% =========================================================================
% Back matter
% =========================================================================
\backmatter

% Declarations: 아래 각 항목을 원고에 맞게 영문으로 작성합니다.
%   Funding(연구비) / Competing interests(이해상충) / Ethics approval(IRB 승인) /
%   Informed consent(연구참여 동의, 미성년자는 보호자 동의) /
%   Data availability(데이터 공개) / Generative AI use(생성형 AI 사용 범위)
\bmhead{Declarations}
\begin{itemize}
  \item \textbf{Funding.} State funding sources, or that no funding was received.
  \item \textbf{Competing interests.} State any competing interests.
  \item \textbf{Ethics approval.} State IRB/ethics approval information.
  \item \textbf{Informed consent.} State that informed consent (and guardian
        consent for minors) was obtained.
  \item \textbf{Data availability.} State how the data can be accessed or requested.
  \item \textbf{Generative AI use.} Disclose whether and how generative AI was used.
\end{itemize}

% 참고문헌(BibTeX, APA via sn-apacite). \bibliographystyle 은 sn-jnl 이 자동 지정.
\bibliography{sub/references-sn}

\end{document}
%%%%%%%%%%%%%%%%%%%%%%%%%%%%%%%%%%%%%%%%%%%%%%%%%%%%%%%%%%%%%%%%%%%%%
%  END RISE_main.tex
%%%%%%%%%%%%%%%%%%%%%%%%%%%%%%%%%%%%%%%%%%%%%%%%%%%%%%%%%%%%%%%%%%%%%
  로 컴파일합니다(제공된 빌드 스크립트의
% `anon` 명령이 자동 처리하여 RISE_main_anon.pdf 를 만듭니다). \ANON 이 정의되면
% 저자·소속·이메일·ORCID·사사(Funding)가 익명화됩니다. 저자 식별 정보는 익명 원고와
% 별도로 title-page.tex(투고용 표지)로 제출합니다.
\newif\ifRISEanon
\ifdefined\ANON\RISEanontrue\fi

% =========================================================================
% Front matter (sn-jnl 양식) — 아래 제목/저자/초록/키워드를 새 원고에 맞게 수정
% =========================================================================
\title[Short Running Title]{Full Article Title Goes Here:
Subtitle If Any}

% 저자 순서대로 \author 로 나열하고, 교신저자는 \author* 로 표시합니다.
% \ifRISEanon 이 참이면(익명본) 저자·소속·이메일이 익명 처리됩니다.
\ifRISEanon
  \author*[1]{\fnm{}\sur{Author name withheld for double-blind review}}
  \affil[1]{\orgname{Affiliation withheld for double-blind review}}
\else
  \author[1]{\fnm{First} \sur{Author}\,\orcidlink{0000-0000-0000-0000}}
  \author*[1]{\fnm{Corresponding} \sur{Author}\,\orcidlink{0000-0000-0000-0000}}\email{corresponding.author@example.edu}
  \author[2]{\fnm{Third} \sur{Author}}
  \affil[1]{\orgname{Your Institution}, \orgaddress{\city{City}, \country{Country}}}
  \affil[2]{\orgname{Second Institution}, \orgaddress{\city{City}, \country{Country}}}
\fi

% 초록: Research in Science Education 은 보통 구조화되지 않은(unstructured) 단일
% 문단 초록을 사용합니다. 배경/문제의식 → 목적·연구질문 → 대상·방법 → 핵심 결과
% → 과학교육 연구·실천에의 함의를 한 문단으로 압축합니다. 최신 저자 안내에서 분량
% 제한을 확인하세요. (이 클래스는 pdfLaTeX 로 컴파일되므로 본문은 영문으로 작성합니다.)
\abstract{Write the abstract here as a single unstructured paragraph. State the
background and problem, the purpose and research questions, the participants and
methods, the key findings, and the implications for science-education research
and practice. Check the journal's current author guidelines for length limits.}

\keywords{keyword one, keyword two, keyword three, keyword four}

\begin{document}

\maketitle

% =========================================================================
% 본문 — 각 절 파일은 sub/ 폴더에 있습니다(예시 내용 없이 안내만 포함).
% =========================================================================
%%%%%%%%%%%%%%%%%%%%%%%%%%%%%%%%%%%%%%%%%%%%%%%%%%%%%%%%%%%%%%%%%%%%%
%  sub/en-1-introduction.tex  —  1. Introduction
%%%%%%%%%%%%%%%%%%%%%%%%%%%%%%%%%%%%%%%%%%%%%%%%%%%%%%%%%%%%%%%%%%%%%
\section{Introduction}\label{sec:introduction}

% 서론에서는 연구가 다루는 과학교육 문제와 그 중요성을 제시하고, 선행연구를
% 비판적으로 검토하여 연구 공백(gap)을 드러낸 뒤, 본 연구의 목적과 연구 질문을
% 명확히 밝힙니다. Research in Science Education 은 이론적 기여와 실증적 근거의
% 연결을 중시하므로, 연구 질문이 어떤 이론적·실천적 함의를 겨냥하는지 초반에
% 분명히 하세요.

Write the introduction here. State the science-education problem, review prior
work to expose the gap, and end with explicit research questions.

% 인용은 \parencite{key} (괄호 인용) 또는 \textcite{key} (본문 인용)로 씁니다.
% 예: Prior work has examined this issue \parencite{example2024}, and
% \textcite{example2024} argued that \dots
   % 1. Introduction
%%%%%%%%%%%%%%%%%%%%%%%%%%%%%%%%%%%%%%%%%%%%%%%%%%%%%%%%%%%%%%%%%%%%%
%  sub/en-2-framework.tex  —  2. Theoretical Framework
%%%%%%%%%%%%%%%%%%%%%%%%%%%%%%%%%%%%%%%%%%%%%%%%%%%%%%%%%%%%%%%%%%%%%
\section{Theoretical Framework}\label{sec:framework}

% 연구를 뒷받침하는 이론적 틀, 핵심 개념, 분석 렌즈를 제시합니다. 연구 질문과
% 방법(코딩 체계, 루브릭 등)을 이 틀에 연결하면 이후 결과 해석의 근거가 됩니다.

Present the theoretical framework, key constructs, and analytic lens that ground
the study and connect the research questions to the methods.
      % 2. Theoretical Framework
%%%%%%%%%%%%%%%%%%%%%%%%%%%%%%%%%%%%%%%%%%%%%%%%%%%%%%%%%%%%%%%%%%%%%
%  sub/en-3-methods.tex  —  3. Methods
%%%%%%%%%%%%%%%%%%%%%%%%%%%%%%%%%%%%%%%%%%%%%%%%%%%%%%%%%%%%%%%%%%%%%
\section{Methods}\label{sec:methods}

% 연구 설계, 맥락과 참여자, 자료 수집, 분석 절차(코딩 체계·신뢰도·타당도)를
% 재현 가능하도록 기술합니다. 연구윤리(IRB) 승인과 동의 절차도 밝힙니다.

Describe the research design, context and participants, data sources, and the
analysis procedure (coding scheme, reliability, validity) in enough detail to be
reproducible.

% 표 예시 (필요 시 사용):
% \begin{table}[htbp]
%   \centering
%   \caption{Participant summary}\label{tab:participants}
%   \begin{tabularx}{\linewidth}{@{}lL@{}}
%     \toprule
%     Group & Description \\
%     \midrule
%     A & \dots \\
%     \bottomrule
%   \end{tabularx}
%   \TableNote{Add a table note here if needed.}
% \end{table}
        % 3. Methods
%%%%%%%%%%%%%%%%%%%%%%%%%%%%%%%%%%%%%%%%%%%%%%%%%%%%%%%%%%%%%%%%%%%%%
%  sub/en-4-results.tex  —  4. Results
%%%%%%%%%%%%%%%%%%%%%%%%%%%%%%%%%%%%%%%%%%%%%%%%%%%%%%%%%%%%%%%%%%%%%
\section{Results}\label{sec:results}

% 연구 질문 순서에 따라 결과를 제시합니다. 표와 그림은 본문에서 먼저 언급한 뒤
% 배치하고, 모든 표/그림에 \label 을 붙여 \ref/\cref 로 참조합니다. 해석은
% 최소화하고 근거(수치·인용·사례)를 중심으로 서술합니다.

Report the findings organized by research question. Reference every table and
figure in the text before it appears.

% 그림 예시 (필요 시 사용):
% \begin{figure}[htbp]
%   \centering
%   \includegraphics[width=\linewidth]{images/your-figure}
%   \caption{Describe the figure.}\label{fig:example}
% \end{figure}
        % 4. Results
%%%%%%%%%%%%%%%%%%%%%%%%%%%%%%%%%%%%%%%%%%%%%%%%%%%%%%%%%%%%%%%%%%%%%
%  sub/en-5-discussion.tex  —  5. Discussion and Conclusion
%%%%%%%%%%%%%%%%%%%%%%%%%%%%%%%%%%%%%%%%%%%%%%%%%%%%%%%%%%%%%%%%%%%%%
\section{Discussion and Conclusion}\label{sec:discussion}

% 핵심 결과를 연구 질문·이론 틀과 연결해 해석하고, 선행연구와의 관계(확장/반박),
% 과학교육 연구·실천에의 함의, 연구의 한계와 후속 연구 방향을 제시합니다.
% Research in Science Education 은 이론적·교육적 기여를 분명히 드러내는 논의를
% 중시합니다.

Interpret the findings in relation to the research questions and framework,
connect to prior work, state implications for science-education research and
practice, and acknowledge limitations and directions for future work.
     % 5. Discussion and Conclusion

% =========================================================================
% Back matter
% =========================================================================
\backmatter

% Declarations: 아래 각 항목을 원고에 맞게 영문으로 작성합니다.
%   Funding(연구비) / Competing interests(이해상충) / Ethics approval(IRB 승인) /
%   Informed consent(연구참여 동의, 미성년자는 보호자 동의) /
%   Data availability(데이터 공개) / Generative AI use(생성형 AI 사용 범위)
\bmhead{Declarations}
\begin{itemize}
  \item \textbf{Funding.} State funding sources, or that no funding was received.
  \item \textbf{Competing interests.} State any competing interests.
  \item \textbf{Ethics approval.} State IRB/ethics approval information.
  \item \textbf{Informed consent.} State that informed consent (and guardian
        consent for minors) was obtained.
  \item \textbf{Data availability.} State how the data can be accessed or requested.
  \item \textbf{Generative AI use.} Disclose whether and how generative AI was used.
\end{itemize}

% 참고문헌(BibTeX, APA via sn-apacite). \bibliographystyle 은 sn-jnl 이 자동 지정.
\bibliography{sub/references-sn}

\end{document}
%%%%%%%%%%%%%%%%%%%%%%%%%%%%%%%%%%%%%%%%%%%%%%%%%%%%%%%%%%%%%%%%%%%%%
%  END RISE_main.tex
%%%%%%%%%%%%%%%%%%%%%%%%%%%%%%%%%%%%%%%%%%%%%%%%%%%%%%%%%%%%%%%%%%%%%
  로 컴파일합니다(제공된 빌드 스크립트의
% `anon` 명령이 자동 처리하여 RISE_main_anon.pdf 를 만듭니다). \ANON 이 정의되면
% 저자·소속·이메일·ORCID·사사(Funding)가 익명화됩니다. 저자 식별 정보는 익명 원고와
% 별도로 title-page.tex(투고용 표지)로 제출합니다.
\newif\ifRISEanon
\ifdefined\ANON\RISEanontrue\fi

% =========================================================================
% Front matter (sn-jnl 양식) — 아래 제목/저자/초록/키워드를 새 원고에 맞게 수정
% =========================================================================
\title[Short Running Title]{Full Article Title Goes Here:
Subtitle If Any}

% 저자 순서대로 \author 로 나열하고, 교신저자는 \author* 로 표시합니다.
% \ifRISEanon 이 참이면(익명본) 저자·소속·이메일이 익명 처리됩니다.
\ifRISEanon
  \author*[1]{\fnm{}\sur{Author name withheld for double-blind review}}
  \affil[1]{\orgname{Affiliation withheld for double-blind review}}
\else
  \author[1]{\fnm{First} \sur{Author}\,\orcidlink{0000-0000-0000-0000}}
  \author*[1]{\fnm{Corresponding} \sur{Author}\,\orcidlink{0000-0000-0000-0000}}\email{corresponding.author@example.edu}
  \author[2]{\fnm{Third} \sur{Author}}
  \affil[1]{\orgname{Your Institution}, \orgaddress{\city{City}, \country{Country}}}
  \affil[2]{\orgname{Second Institution}, \orgaddress{\city{City}, \country{Country}}}
\fi

% 초록: Research in Science Education 은 보통 구조화되지 않은(unstructured) 단일
% 문단 초록을 사용합니다. 배경/문제의식 → 목적·연구질문 → 대상·방법 → 핵심 결과
% → 과학교육 연구·실천에의 함의를 한 문단으로 압축합니다. 최신 저자 안내에서 분량
% 제한을 확인하세요. (이 클래스는 pdfLaTeX 로 컴파일되므로 본문은 영문으로 작성합니다.)
\abstract{Write the abstract here as a single unstructured paragraph. State the
background and problem, the purpose and research questions, the participants and
methods, the key findings, and the implications for science-education research
and practice. Check the journal's current author guidelines for length limits.}

\keywords{keyword one, keyword two, keyword three, keyword four}

\begin{document}

\maketitle

% =========================================================================
% 본문 — 각 절 파일은 sub/ 폴더에 있습니다(예시 내용 없이 안내만 포함).
% =========================================================================
%%%%%%%%%%%%%%%%%%%%%%%%%%%%%%%%%%%%%%%%%%%%%%%%%%%%%%%%%%%%%%%%%%%%%
%  sub/en-1-introduction.tex  —  1. Introduction
%%%%%%%%%%%%%%%%%%%%%%%%%%%%%%%%%%%%%%%%%%%%%%%%%%%%%%%%%%%%%%%%%%%%%
\section{Introduction}\label{sec:introduction}

% 서론에서는 연구가 다루는 과학교육 문제와 그 중요성을 제시하고, 선행연구를
% 비판적으로 검토하여 연구 공백(gap)을 드러낸 뒤, 본 연구의 목적과 연구 질문을
% 명확히 밝힙니다. Research in Science Education 은 이론적 기여와 실증적 근거의
% 연결을 중시하므로, 연구 질문이 어떤 이론적·실천적 함의를 겨냥하는지 초반에
% 분명히 하세요.

Write the introduction here. State the science-education problem, review prior
work to expose the gap, and end with explicit research questions.

% 인용은 \parencite{key} (괄호 인용) 또는 \textcite{key} (본문 인용)로 씁니다.
% 예: Prior work has examined this issue \parencite{example2024}, and
% \textcite{example2024} argued that \dots
   % 1. Introduction
%%%%%%%%%%%%%%%%%%%%%%%%%%%%%%%%%%%%%%%%%%%%%%%%%%%%%%%%%%%%%%%%%%%%%
%  sub/en-2-framework.tex  —  2. Theoretical Framework
%%%%%%%%%%%%%%%%%%%%%%%%%%%%%%%%%%%%%%%%%%%%%%%%%%%%%%%%%%%%%%%%%%%%%
\section{Theoretical Framework}\label{sec:framework}

% 연구를 뒷받침하는 이론적 틀, 핵심 개념, 분석 렌즈를 제시합니다. 연구 질문과
% 방법(코딩 체계, 루브릭 등)을 이 틀에 연결하면 이후 결과 해석의 근거가 됩니다.

Present the theoretical framework, key constructs, and analytic lens that ground
the study and connect the research questions to the methods.
      % 2. Theoretical Framework
%%%%%%%%%%%%%%%%%%%%%%%%%%%%%%%%%%%%%%%%%%%%%%%%%%%%%%%%%%%%%%%%%%%%%
%  sub/en-3-methods.tex  —  3. Methods
%%%%%%%%%%%%%%%%%%%%%%%%%%%%%%%%%%%%%%%%%%%%%%%%%%%%%%%%%%%%%%%%%%%%%
\section{Methods}\label{sec:methods}

% 연구 설계, 맥락과 참여자, 자료 수집, 분석 절차(코딩 체계·신뢰도·타당도)를
% 재현 가능하도록 기술합니다. 연구윤리(IRB) 승인과 동의 절차도 밝힙니다.

Describe the research design, context and participants, data sources, and the
analysis procedure (coding scheme, reliability, validity) in enough detail to be
reproducible.

% 표 예시 (필요 시 사용):
% \begin{table}[htbp]
%   \centering
%   \caption{Participant summary}\label{tab:participants}
%   \begin{tabularx}{\linewidth}{@{}lL@{}}
%     \toprule
%     Group & Description \\
%     \midrule
%     A & \dots \\
%     \bottomrule
%   \end{tabularx}
%   \TableNote{Add a table note here if needed.}
% \end{table}
        % 3. Methods
%%%%%%%%%%%%%%%%%%%%%%%%%%%%%%%%%%%%%%%%%%%%%%%%%%%%%%%%%%%%%%%%%%%%%
%  sub/en-4-results.tex  —  4. Results
%%%%%%%%%%%%%%%%%%%%%%%%%%%%%%%%%%%%%%%%%%%%%%%%%%%%%%%%%%%%%%%%%%%%%
\section{Results}\label{sec:results}

% 연구 질문 순서에 따라 결과를 제시합니다. 표와 그림은 본문에서 먼저 언급한 뒤
% 배치하고, 모든 표/그림에 \label 을 붙여 \ref/\cref 로 참조합니다. 해석은
% 최소화하고 근거(수치·인용·사례)를 중심으로 서술합니다.

Report the findings organized by research question. Reference every table and
figure in the text before it appears.

% 그림 예시 (필요 시 사용):
% \begin{figure}[htbp]
%   \centering
%   \includegraphics[width=\linewidth]{images/your-figure}
%   \caption{Describe the figure.}\label{fig:example}
% \end{figure}
        % 4. Results
%%%%%%%%%%%%%%%%%%%%%%%%%%%%%%%%%%%%%%%%%%%%%%%%%%%%%%%%%%%%%%%%%%%%%
%  sub/en-5-discussion.tex  —  5. Discussion and Conclusion
%%%%%%%%%%%%%%%%%%%%%%%%%%%%%%%%%%%%%%%%%%%%%%%%%%%%%%%%%%%%%%%%%%%%%
\section{Discussion and Conclusion}\label{sec:discussion}

% 핵심 결과를 연구 질문·이론 틀과 연결해 해석하고, 선행연구와의 관계(확장/반박),
% 과학교육 연구·실천에의 함의, 연구의 한계와 후속 연구 방향을 제시합니다.
% Research in Science Education 은 이론적·교육적 기여를 분명히 드러내는 논의를
% 중시합니다.

Interpret the findings in relation to the research questions and framework,
connect to prior work, state implications for science-education research and
practice, and acknowledge limitations and directions for future work.
     % 5. Discussion and Conclusion

% =========================================================================
% Back matter
% =========================================================================
\backmatter

% Declarations: 아래 각 항목을 원고에 맞게 영문으로 작성합니다.
%   Funding(연구비) / Competing interests(이해상충) / Ethics approval(IRB 승인) /
%   Informed consent(연구참여 동의, 미성년자는 보호자 동의) /
%   Data availability(데이터 공개) / Generative AI use(생성형 AI 사용 범위)
\bmhead{Declarations}
\begin{itemize}
  \item \textbf{Funding.} State funding sources, or that no funding was received.
  \item \textbf{Competing interests.} State any competing interests.
  \item \textbf{Ethics approval.} State IRB/ethics approval information.
  \item \textbf{Informed consent.} State that informed consent (and guardian
        consent for minors) was obtained.
  \item \textbf{Data availability.} State how the data can be accessed or requested.
  \item \textbf{Generative AI use.} Disclose whether and how generative AI was used.
\end{itemize}

% 참고문헌(BibTeX, APA via sn-apacite). \bibliographystyle 은 sn-jnl 이 자동 지정.
\bibliography{sub/references-sn}

\end{document}
%%%%%%%%%%%%%%%%%%%%%%%%%%%%%%%%%%%%%%%%%%%%%%%%%%%%%%%%%%%%%%%%%%%%%
%  END RISE_main.tex
%%%%%%%%%%%%%%%%%%%%%%%%%%%%%%%%%%%%%%%%%%%%%%%%%%%%%%%%%%%%%%%%%%%%%
